\documentclass[conference]{IEEEtran}
\IEEEoverridecommandlockouts
\usepackage{cite}
\usepackage{amsmath,amssymb,amsfonts}
\usepackage{algorithmic}
\usepackage{graphicx}
\usepackage{textcomp}
\usepackage{xcolor}
\usepackage{booktabs}
\usepackage{multirow}
\usepackage{stfloats}
\usepackage{float}
\usepackage{hyperref}
\def\BibTeX{{\rm B\kern-.05em{\sc i\kern-.025em b}\kern-.08em
    T\kern-.1667em\lower.7ex\hbox{E}\kern-.125emX}}

% ---------------------------------------------------------------
%  Auto-generated dataset macros (filled by MetaStatsAnalysis)
% ---------------------------------------------------------------
\newcommand{\numTrades}{72{,}134{,}741}
\newcommand{\numTradesMillions}{72.1}
\newcommand{\totalVolumeBillions}{18.3}
\newcommand{\numMarkets}{7{,}314{,}375}
\newcommand{\numEvents}{1{,}197{,}300}

\begin{document}

% ---------------------------------------------------------------
%  TITLE
% ---------------------------------------------------------------
\title{Who Profits in Binary Prediction Markets? Maker-Taker Dynamics, Behavioral Bias,
       and Sentiment Arbitrage on Kalshi}

\author{%
  \IEEEauthorblockN{Abhinav Gupta}
  \IEEEauthorblockA{%
    \textit{IIT Roorkee}\\
    Roorkee, India\\
    abhinav\_g@ma.iitr.ac.in}
}

\maketitle

% ---------------------------------------------------------------
%  ABSTRACT
% ---------------------------------------------------------------
\begin{abstract}
Whether liquidity providers in binary prediction markets systematically profit
from uninformed order flow, and whether behavioral biases sustain that profit,
remain open empirical questions. We address them using \numTradesMillions{}~million
executed trades (\$\totalVolumeBillions{}B notional) on Kalshi from July 2021
through November 2025, supplemented by 404.5~million on-chain trades from
Polymarket, applying volume-weighted return decomposition, calibration analysis,
OLS regression, and paired $t$-tests.

A statistically significant Favorite-Longshot Bias (FLSB) emerges on both
platforms, but with a cross-platform contrast: Kalshi longshots are overpriced
($\delta_b < 0$, $p < 0.001$ for all sub-25\% buckets) while Polymarket
favorites above 55\% implied probability are systematically underpriced
($\delta_b > 0$, $t$ up to 361), implicating market design as the key
moderating factor. On the supply side, makers earn $+0.83$~percentage points~(pp)
in volume-weighted returns ($p = 0.0055$), constituting a \$142.5M wealth
transfer that exceeds estimated bid-ask spread compensation by $0.13$\,pp.
This alpha is not uniform: it ranges from $+0.29$\,pp in Finance to
$+3.01$\,pp in World Events, with high-volume categories spanning $+0.82$\,pp
(Sports) to $+1.78$\,pp (Media), consistent with adverse selection varying
by information environment.
On the demand side, takers allocate 63.8\% of capital to YES contracts despite
earning $0.32$\,pp less than NO-side takers, an affirmative framing bias
unexplained by differential risk. Contrary to efficiency-convergence predictions,
maker alpha trends upward at $+8.21$\,bp/month ($R^2 = 0.70$) as retail
participation grows. A sport-level decomposition of the dominant sports category
(74\% of volume) reveals maker alpha ranging from $+0.18$\,pp in MLB to
$+3.43$\,pp in UFC/Boxing, with no significant time-to-resolution decay
($p = 0.88$), providing direct evidence of persistent fan-driven sentiment
arbitrage. Taken together, these findings indicate that binary prediction markets
structurally disadvantage uninformed retail participants and call for targeted
disclosure requirements and calibration-oriented regulatory oversight.
\end{abstract}

\begin{IEEEkeywords}
prediction markets, market microstructure, favorite-longshot bias,
maker-taker dynamics, behavioral finance, information efficiency, Kalshi
\end{IEEEkeywords}

% ---------------------------------------------------------------
%  I. PREDICTION MARKETS AND KALSHI
% ---------------------------------------------------------------
\section{Prediction Markets and Kalshi}
\label{sec:intro}

Binary prediction markets occupy a conceptual middle ground between financial
derivatives and information-aggregation mechanisms \cite{wolfers2004,arrow2008,hanson2003,berg2008}.
Participants stake capital on discrete real-world outcomes (election results,
economic releases, sporting events) and receive a binary payoff on resolution.
The market price $P_t \in (0,1)$ represents the collective risk-neutral
probability assigned to the event at time $t$.

Kalshi, the first CFTC-regulated binary prediction market in the United States,
operates a central limit order book (CLOB) in which participants submit resting
limit orders or market orders. For each binary contract, a YES position pays
\$1 if the event resolves positively and \$0 otherwise; the complementary NO
position pays \$1 if it does not. A trade at price $P_t$ implies YES cost
basis $= P_t$ and NO cost basis $= 1 - P_t$.

Despite Kalshi's growing regulatory prominence, the microstructure of binary
prediction markets remains understudied. Prior work has focused on price
calibration \cite{wolfers2004} and information aggregation \cite{arrow2008},
but systematic evidence on maker-taker dynamics, behavioral biases, and their
interaction at the trade level is absent. We fill this gap using
\numTradesMillions{} million trades across \numMarkets{} Kalshi markets
spanning July 2021 through November 2025.

\textbf{Contributions.} We make six contributions:
\begin{enumerate}
    \item We document a robust Favorite-Longshot Bias on Kalshi and a
          complementary reverse-favorite bias on Polymarket, identifying
          market design as the key moderator of calibration failure.
    \item We quantify a \$142.5M structural wealth transfer from takers to
          makers, showing that maker alpha ($+0.83$\,pp) significantly exceeds
          the estimated bid-ask spread (half-spread $= 0.70$\,pp, $t = 84.6$).
    \item We identify a systematic YES-side affirmative bias in taker order
          flow (63.8\% capital share), with YES-side takers earning
          significantly lower returns than NO-side takers.
    \item We show that maker alpha varies by category consistent with adverse
          selection theory \cite{kyle1985}: highest in sentiment-driven markets
          (World Events $+3.01$\,pp, Media $+1.78$\,pp, Entertainment $+1.37$\,pp)
          and lowest among major categories where takers are informed
          (Finance, $+0.29$\,pp); Esports, a minor category, exhibits negative maker alpha.
    \item We document a statistically significant upward trend in maker alpha
          ($\hat{\gamma}_1 = +8.21$\,bp/month, $R^2 = 0.70$), inconsistent
          with standard efficiency-convergence predictions \cite{fama1970}.
    \item We introduce the concept of sports sentiment arbitrage, decomposing
          sports markets (74\% of total volume) into ten sports and finding
          universal YES-side fan bias ($>72\%$) and sport-specific maker alpha
          from $+0.18$\,pp (MLB) to $+3.43$\,pp (UFC/Boxing), with no
          significant time-to-resolution decay ($p = 0.88$).
\end{enumerate}

The remainder of this paper is organized as follows. Section~\ref{sec:data}
describes data and methodology. Sections~\ref{sec:flsb}--\ref{sec:behavioral}
present the empirical findings. Section~\ref{sec:discussion} provides
a unified interpretation. Section~\ref{sec:limitations} discusses limitations
and Section~\ref{sec:conclusion} concludes.

% ---------------------------------------------------------------
%  II. DATA AND METHODOLOGY
% ---------------------------------------------------------------
\section{Data and Methodology}
\label{sec:data}

\subsection{Binary Contract Payoff Structure}

Let $Y_i \in \{0, 1\}$ denote the terminal payoff of contract $i$ at expiration $T$.
The equilibrium market price $P_t \in (0,1)$ represents the aggregate risk-neutral
subjective probability:
\begin{equation}
  P_t = \mathbb{E}_t[Y_i].
  \label{eq:price}
\end{equation}

Trades are classified by order-submission type. A \emph{taker} submits a
market order that immediately executes against a resting limit order; a
\emph{maker} provides resting liquidity by posting limit orders. For each trade:
\begin{align}
  \Pi_{\text{Taker}} &=
  \begin{cases}
    Y_i - P_t          & \text{if taker buys YES at } P_t, \\
    (1-Y_i) - (1-P_t)  & \text{if taker buys NO at } (1-P_t),
  \end{cases} \label{eq:pnl}\\
  \Pi_{\text{Maker}}  &= -\Pi_{\text{Taker}}, \label{eq:maker_pnl}
\end{align}
since prediction markets are zero-sum (ignoring fees) \cite{manski2006}.

We define the volume-weighted average return (VWAR):
\begin{equation}
  \overline{R} = \frac{\sum_j \Pi_j \cdot n_j}{\sum_j n_j},
  \label{eq:vwar}
\end{equation}
where $n_j$ is the contract count for trade $j$.

\subsection{Sample}
\label{sec:sample}

Our primary dataset covers Kalshi trades from platform launch through Q4 2025,
comprising approximately \numTradesMillions{} million executed trades across
\numMarkets{} finalized markets and \numEvents{} distinct events.
Cross-platform evidence is drawn from Polymarket's on-chain Conditional Token
Framework (CTF) exchange, covering 404.5~million trades across 383,707 resolved
markets from September 2020 through February 2026.
Polymarket markets resolved YES in 40.8\% of cases versus 23.9\% on Kalshi,
reflecting Kalshi's prevalence of low-probability Sports and Crypto contracts.
Because Polymarket operates a hybrid AMM/CLOB system, the maker-taker PnL
framework applies only to Kalshi; Polymarket provides cross-platform calibration
and volume-concentration benchmarks.

Table~\ref{tab:market_overview} summarises Kalshi market coverage by category group.

\begin{table}[htbp]
  \caption{Market Architecture Summary by Category (Kalshi, 2021--2025)}
  \label{tab:market_overview}
  \begin{center}
  \begin{tabular}{lrrrr}
    \toprule
    \textbf{Category} & \textbf{N} & \textbf{\%YES} & \textbf{Med.\ Dur.\ (h)} & \textbf{Mean Dur.\ (h)} \\
    \midrule
    Sports     & 3,532,947 & 11.2 & 14 & 27.3 \\
    Crypto     & 1,788,112 & 26.3 &  1 &  3.2 \\
    Finance    & 1,612,438 & 49.9 &  1 & 19.9 \\
    Politics   &   123,966 & 33.6 & 24 & 100.0 \\
    Other      &   256,912 & 13.5 &  8 & 117.1 \\
    \midrule
    \textbf{All} & \numMarkets{} & 23.9 & 7 & 24.2 \\
    \bottomrule
  \end{tabular}
  \end{center}
\end{table}

\begin{figure}[H]
  \centerline{\includegraphics[width=\columnwidth]{output/cp01_market_overview/kalshi/kalshi_cp01_market_overview.pdf}}
  \caption{Kalshi market overview: trade volume, market count, and resolution rates
           by category.}
  \label{fig:market_overview}
\end{figure}

\begin{figure}[H]
  \centerline{\includegraphics[width=\columnwidth]{output/cp01_market_overview/polymarket/polymarket_cp01_market_overview.pdf}}
  \caption{Polymarket market overview: monthly CTF trade volume and resolution funnel.}
  \label{fig:market_overview_poly}
\end{figure}

% ---------------------------------------------------------------
%  III. THE LONGSHOT BIAS ON KALSHI
% ---------------------------------------------------------------
\section{The Longshot Bias on Kalshi}
\label{sec:flsb}

Group contracts into price buckets $b \in (0,1)$ based on transaction price $P_t$.
Let $f(b)$ be the empirical win frequency of taker-YES trades within bucket $b$,
and $\bar{P}_b$ the average implied probability in $b$. The calibration deviation is:
\begin{equation}
  \delta_b = f(b) - \bar{P}_b.
  \label{eq:delta_b}
\end{equation}
The Favorite-Longshot Bias (FLSB) predicts $\delta_b < 0$ for small $\bar{P}_b$
(longshots overpriced) and $\delta_b > 0$ for large $\bar{P}_b$ (favorites
underpriced) \cite{ali1977,thaler1988,snowberg2010,camerer1998,ottaviani2008}.

Figure~\ref{fig:flsb} shows the calibration curve for Kalshi. We group trades
into 5-cent bins and compute $\delta_b$ for each. For buckets with
$\bar{P}_b < 0.25$, a one-sample $t$-test rejects the null $\delta_b = 0$
($t < -16.3$ across all sub-25\% buckets, $p < 0.001$), confirming a
statistically significant FLSB on Kalshi.

Figure~\ref{fig:flsb_poly} presents the Polymarket calibration curve, which
reveals a qualitatively distinct and complementary pattern. Longshots
($\bar{P}_b < 0.25$) are again overpriced ($\delta_b < 0$, all $p < 0.001$),
consistent with the Kalshi finding. However, Polymarket exhibits a pronounced
\emph{reverse} bias at high implied probabilities: favorites with
$\bar{P}_b > 0.55$ are systematically \emph{underpriced} ($\delta_b > 0$),
with $t$-statistics reaching 297 at the 60--65\% bucket and 361 at the
95--97.5\% bucket. Takers who purchase near-certain contracts on Polymarket
earn statistically significant \emph{positive} returns. This pattern is
absent on Kalshi and likely reflects Polymarket's AMM-driven pricing
mechanism, which mechanically underestimates the probability of near-certain
events as liquidity thins near the boundary. Together, the two platforms
confirm that calibration failure is a structural feature of binary prediction
markets, with its direction modulated by market design.

\begin{figure}[H]
  \centerline{\includegraphics[width=\columnwidth]{output/cp02_longshot_bias/kalshi/kalshi_cp02_longshot_bias.pdf}}
  \caption{Kalshi calibration curve and $\delta_b$ by price bucket.
           Negative bars at low prices confirm the Favorite-Longshot Bias.}
  \label{fig:flsb}
\end{figure}

\begin{figure}[H]
  \centerline{\includegraphics[width=\columnwidth]{output/cp02_longshot_bias/polymarket/polymarket_cp02_longshot_bias.pdf}}
  \caption{Polymarket calibration deviation $\delta_b$: cross-platform replication reveals reverse bias at high implied probabilities.}
  \label{fig:flsb_poly}
\end{figure}

% ---------------------------------------------------------------
%  IV. THE MAKER-TAKER WEALTH TRANSFER
% ---------------------------------------------------------------
\section{The Maker-Taker Wealth Transfer}
\label{sec:maker_taker}

Table~\ref{tab:maker_taker} presents aggregate maker and taker VWARs.
Makers earn a positive aggregate VWAR of \textbf{0.83} pp while takers
earn an equivalent negative return; the cumulative wealth transfer is
approximately \textbf{\$142.5M}. A monthly t-test (one sample, equally
weighted across calendar months) yields $t = -2.90$ ($p = 0.0055$).
The negative sign reflects that the unweighted monthly mean is depressed
by thin early-platform quarters (2021--2022); the positive aggregate
VWAR is driven by high-volume later periods in which maker alpha is
substantially larger.

\begin{table}[htbp]
  \caption{Maker vs.\ Taker Volume-Weighted Average Returns}
  \label{tab:maker_taker}
  \begin{center}
  \begin{tabular}{lrr}
    \toprule
    & \textbf{Maker} & \textbf{Taker} \\
    \midrule
    VWAR (pp)       & $+0.83$ & $-0.83$ \\
    Total PnL (\$)  & $+142.5$M & $-142.5$M \\
    N trades (finalized) & \multicolumn{2}{c}{$67{,}761{,}406$} \\
    $t$-stat        & $-2.90$ & $+2.90$ \\
    $p$-value       & \multicolumn{2}{c}{$0.0055$} \\
    \bottomrule
  \end{tabular}
  \end{center}
\end{table}

\begin{figure}[H]
  \centerline{\includegraphics[width=\columnwidth]{output/cp03_maker_taker_pnl/kalshi/kalshi_cp03_maker_taker_pnl.pdf}}
  \caption{Maker vs.\ Taker volume-weighted returns by price bucket.
           Green bars show maker alpha; red bars show taker losses.}
  \label{fig:maker_taker}
\end{figure}

\subsection{Decomposing Returns by Role}
\label{sec:decompose}

Figure~\ref{fig:temporal_dual} decomposes quarterly returns by role from market
inception (July 2021) through Q4 2025. Maker VWR (blue) trends upward from negative values in the earliest quarters
(thin market, dominated by sophisticated early adopters) to persistently positive
alpha by 2023--2025 as retail taker flow expanded. Taker VWR (green) mirrors
this trajectory with reversed sign. The OLS time-trend ($\hat{\gamma}_1 = +8.21$
bp/month, $R^2 = 0.70$) indicates that growing platform scale has
\emph{amplified} rather than compressed the structural wealth transfer.

\begin{figure*}[!t]
  \centerline{\includegraphics[width=\textwidth]{output/cp06_temporal_dynamics/kalshi/kalshi_cp06_temporal_dynamics.pdf}}
  \caption{Temporal dynamics of maker-taker returns on Kalshi (July 2021--November 2025).
           \textit{Top}: Quarterly maker (blue) and taker (green) VWR over time.
           \textit{Bottom-left}: Taker volume composition by price bucket (1--10¢ through 91--99¢)
           as a stacked area per quarter.
           \textit{Bottom-right}: Monthly maker VWR (bubbles, sized by volume) with OLS
           time-trend $\hat{\gamma}_1$ and 95\% confidence band.}
  \label{fig:temporal_dual}
\end{figure*}

% ---------------------------------------------------------------
%  V. IS THIS JUST SPREAD COMPENSATION?
% ---------------------------------------------------------------
\section{Is This Just Spread Compensation?}
\label{sec:spread}

Following the realized-spread decomposition of \cite{glosten1985,roll1984,ohara1995,amihud2002},
maker alpha can be written as:
\begin{equation}
  \alpha_{\text{Maker}} = \tfrac{1}{2}S_t - \Delta P_{t+\Delta t},
  \label{eq:alpha}
\end{equation}
where $\tfrac{1}{2}S_t$ is the half-spread captured on liquidity provision and
$\Delta P_{t+\Delta t}$ is the adverse price impact. The effective half-spread
is proxied by half the average absolute price change between consecutive trades
within the same market (estimated at $0.70$\,pp). We test
$H_0: \bar{R}_{\text{Maker}} \leq \hat{S}_t / 2$ against
$H_1: \bar{R}_{\text{Maker}} > \hat{S}_t / 2$.

Figure~\ref{fig:spread} plots the 30-day rolling maker VWAR against the
estimated half-spread. The one-sided $t$-test strongly rejects the null
($t = 84.6$, $p < 0.001$): maker VWAR of $+0.83$\,pp exceeds the half-spread
of $0.70$\,pp, with the residual $+0.13$\,pp representing pure adverse-selection
alpha that cannot be attributed to mechanical spread capture.

\begin{figure}[H]
  \centerline{\includegraphics[width=\columnwidth]{output/cp04_spread_analysis/kalshi/kalshi_cp04_spread_analysis.pdf}}
  \caption{30-day rolling maker VWAR vs.\ estimated half-spread.
           Maker alpha persistently exceeds spread compensation.}
  \label{fig:spread}
\end{figure}

% ---------------------------------------------------------------
%  VI. VARIATION ACROSS CATEGORIES
% ---------------------------------------------------------------
\section{Variation Across Categories}
\label{sec:categories}

Let $g \in \{\text{Politics, Finance, Sports, Crypto, Entertainment, \ldots}\}$
denote the category group. We estimate category-specific maker alpha via OLS:
\begin{equation}
  R_j = \beta_0 + \sum_{g \neq g_0} \beta_g D_{j,g} + \varepsilon_j,
  \label{eq:ols_cat}
\end{equation}
where $D_{j,g} = \mathbf{1}[\text{trade } j \in \text{group } g]$ and
$g_0$ is the baseline (most-traded) category. The $\hat{\beta}_g$ coefficients
test whether each category's maker alpha differs significantly from the Sports
baseline; Figure~\ref{fig:categories} reports the corresponding
volume-weighted average returns (VWARs).

\begin{figure}[H]
  \centerline{\includegraphics[width=\columnwidth]{output/cp05_category_heterogeneity/kalshi/kalshi_cp05_category_heterogeneity.pdf}}
  \caption{Maker vs.\ Taker volume-weighted return by market category, sorted by total volume.
           Green bars: maker VWR; blue bars: taker VWR.}
  \label{fig:categories}
\end{figure}

Figure~\ref{fig:categories} shows the grouped bar chart of maker (green) and
taker (blue) VWR per category. Among major categories, Finance markets exhibit
the lowest maker alpha ($+0.29$\,pp), consistent with informed takers and
tighter adverse selection costs. Esports, a minor category representing
approximately 0.1\% of total volume (22.4M contracts), exhibits negative maker
alpha ($-2.63$\,pp), consistent with a concentrated base of sophisticated,
informed participants. Crypto ($+0.85$\,pp) and Sports ($+0.82$\,pp) cluster near the
cross-market average. World Events ($+3.01$\,pp) registers the highest maker
alpha of all categories; Media ($+1.78$\,pp) and Entertainment ($+1.37$\,pp)
follow, all consistent with sentiment-driven, uninformed order flow. Taker VWR mirrors maker VWR with reversed sign, confirming the
zero-sum structure of the market.

% ---------------------------------------------------------------
%  VII. EVOLUTION OVER TIME
% ---------------------------------------------------------------
\section{Evolution Over Time}
\label{sec:temporal}

Standard market efficiency theory \cite{fama1970} predicts that growing
competition among liquidity providers should compress maker alpha over time.
We test this prediction by regressing the monthly volume-weighted maker return
$R_{M,t}$ on a time index $\tau_t$ (months since platform launch):
\begin{equation}
  R_{M,t} = \gamma_0 + \gamma_1 \tau_t + \varepsilon_t,
  \label{eq:temporal}
\end{equation}
A negative $\hat{\gamma}_1$ would indicate efficiency convergence; a positive
$\hat{\gamma}_1$ would indicate expanding alpha driven by participant-mix effects.

The bottom-right panel of Figure~\ref{fig:temporal_dual} shows monthly maker VWAR
with the OLS trend line and 95\% confidence band. The estimated coefficient
$\hat{\gamma}_1 = +8.21$ bp/month ($p < 0.001$, $R^2 = 0.70$) indicates a
statistically significant \emph{upward} trend: maker alpha has grown over the
sample period, driven by rising retail taker participation rather than
efficiency convergence through Q4 2025.

The stacked area panel (bottom-left of Figure~\ref{fig:temporal_dual}) shows
that taker volume composition by price bucket has remained broadly stable over
the sample period, ruling out a compositional shift toward low-probability
contracts as an explanation for the evolving return gap.

On Polymarket, the monthly CTF trade count grew at a statistically significant
rate over 65 months ($p < 0.001$), confirming that retail participant inflow
was not unique to Kalshi and providing the supply-side condition for expanding
maker alpha across both platforms.

% ---------------------------------------------------------------
%  VIII. THE YES/NO ASYMMETRY
% ---------------------------------------------------------------
\section{The YES/NO Asymmetry}
\label{sec:behavioral}

\subsection{The Asymmetry at Equivalent Prices}
\label{sec:symmetric}

To isolate the directional framing effect from price-level differences, we
compare returns at structurally identical implied probabilities: taker buying
YES at $p$ versus buying NO at implied probability $p$ (i.e., NO price $= 1-P_t = p$).
Under full rationality, returns should be identical at each $p$.

Figure~\ref{fig:symmetric} shows persistent YES underperformance relative to NO
at equivalent price buckets. The top panel of Figure~\ref{fig:symmetric}
decomposes the maker's perspective by position direction. When makers post YES
liquidity (``Maker Bought YES''), their excess return profile differs systematically
from when they post NO liquidity (``Maker Bought NO''), consistent with informed
directional positioning by professional liquidity providers.

\begin{figure*}[!t]
  \centerline{\includegraphics[width=\textwidth]{output/cp08_symmetric_price_asymmetry/kalshi/kalshi_cp08_symmetric_price_asymmetry.pdf}}
  \caption{Symmetric price asymmetry: maker and taker returns decomposed by contract direction.
           \textit{Top}: Maker excess returns at 1-cent granularity, by whether the maker
           holds a YES or NO position.
           \textit{Bottom-left}: Taker VWR for YES@$p$ vs.\ NO@$p$ (grouped bars).
           \textit{Bottom-right}: YES$-$NO taker return gap scatter.}
  \label{fig:symmetric}
\end{figure*}

\subsection{Takers Prefer Affirmative Bets}
\label{sec:yes_no}

We test whether takers systematically prefer YES contracts over economically
equivalent NO contracts. A paired monthly $t$-test on YES-side versus NO-side
returns yields $t = -2.66$ ($p = 0.010$): YES-side takers earn $-0.92$\,pp
versus $-0.61$\,pp for NO-side takers, a statistically significant $-0.32$\,pp
gap. We also measure the affirmative capital tilt via:
\begin{equation}
  \rho_{\text{YES}} = \frac{\sum_j V_j^{\text{YES}}}{\sum_j V_j^{\text{YES}} + \sum_j V_j^{\text{NO}}}.
  \label{eq:vol_ratio}
\end{equation}
We find $\hat{\rho}_{\text{YES}} = 63.8\%$ (binomial $p < 0.001$).
Under symmetric risk, rational participants should split capital equally;
deviation to 63.8\% confirms affirmative framing as a psychological, not
financial, motive \cite{kahneman1979,hong2007}.

\begin{figure}[H]
  \centerline{\includegraphics[width=\columnwidth]{output/cp07_yes_no_asymmetry/kalshi/kalshi_cp07_yes_no_asymmetry.pdf}}
  \caption{YES vs.\ NO taker return and volume split.
           Takers overallocate capital to YES contracts despite earning lower returns than NO-side takers.}
  \label{fig:yes_no}
\end{figure}

\begin{figure}[H]
  \centerline{\includegraphics[width=\columnwidth]{output/cp09_volume_demand_elasticity/kalshi/kalshi_cp09_volume_demand_elasticity.pdf}}
  \caption{YES capital share by market category, with global ratio and binomial test against the 50\% null.}
  \label{fig:volume}
\end{figure}

% ---------------------------------------------------------------
%  IX. DISCUSSION
% ---------------------------------------------------------------
\section{Discussion}
\label{sec:discussion}

\subsection{The Mechanism of Extraction}
\label{sec:mechanism}

The empirical patterns documented above are consistent with a two-tier market
structure: a small set of informed, professional liquidity providers who
systematically exploit the behavioral biases of a larger population of
uninformed, retail-oriented takers. The mechanism operates through three
reinforcing channels:

\textbf{Calibration extraction.} Takers systematically overprice low-probability
events (the FLSB). Makers post resting limit orders at inflated prices, collecting
the miscalibration premium on resolution.

\textbf{Directional framing capture.} Takers display a robust YES-side preference
independent of price level (Section~\ref{sec:yes_no}). Makers accommodate this
demand imbalance by providing NO liquidity at favorable prices, earning structural
alpha through position direction selection (Figure~\ref{fig:symmetric}).

\textbf{Adverse selection gradient.} Maker alpha is largest in categories where
takers are most likely to be sentiment-driven (World Events $+3.01$\,pp,
Media $+1.78$\,pp, Entertainment $+1.37$\,pp) and smallest where takers are
most likely to possess genuine information (Finance $+0.29$\,pp). This gradient
(Section~\ref{sec:categories}) is the hallmark of adverse selection: informed
takers narrow the maker edge.

\subsection{The Professionalization of Liquidity}
\label{sec:hhi}

We measure volume concentration using the Herfindahl-Hirschman Index (HHI)
at the individual contract (ticker) level:
\begin{equation}
  \text{HHI} = \sum_{m=1}^{M} s_m^2,
  \label{eq:hhi}
\end{equation}
where $s_m$ is the volume share of ticker $m$. High concentration indicates
dominance by a small number of high-liquidity contracts. Figure~\ref{fig:hhi}
shows the quarterly HHI trend on Kalshi; Figure~\ref{fig:hhi_poly} shows the
corresponding analysis for Polymarket.

Both platforms report near-zero aggregate HHIs (Kalshi: 0.0008; Polymarket:
0.0013), reflecting the large number of distinct markets in each
($\approx$586,000 and 591,000, respectively). Yet both Lorenz curves reveal
substantial volume concentration hidden by the aggregate metric. On Kalshi,
the top 1\% of tickers account for 76.5\% of all traded volume and the top
10\% account for 94.1\%. On Polymarket, the corresponding figures are 60.5\%
and 87.6\%, consistent with the reported HHI of 0.0013. Across both platforms,
a small subset of high-liquidity contracts dominates total volume; the near-zero
HHI reflects market breadth, not genuine dispersion. The Lorenz-based Gini
coefficient is the more informative concentration measure.

\begin{figure}[H]
  \centerline{\includegraphics[width=\columnwidth]{output/cp10_hhi_concentration/kalshi/kalshi_cp10_hhi_concentration.pdf}}
  \caption{Ticker-level HHI and Lorenz concentration curve, Kalshi.
           Amber zone: moderate concentration (0.15--0.25); red zone: high concentration ($>$0.25).
           Gini coefficient measures inequality of volume distribution across tickers.}
  \label{fig:hhi}
\end{figure}

\begin{figure}[H]
  \centerline{\includegraphics[width=\columnwidth]{output/cp10_hhi_concentration/polymarket/polymarket_cp10_hhi_concentration.pdf}}
  \caption{Polymarket token-level HHI trend and Lorenz concentration curve.
           Despite near-zero aggregate HHI, the top 1\% of tokens account for 60.5\% of volume.}
  \label{fig:hhi_poly}
\end{figure}

\subsection{Category Differences and Participant Selection}
\label{sec:adverse}

We classify market categories by information structure: \emph{scheduled}
(Finance, Macro: CPI releases, FOMC decisions) versus \emph{diffuse}
(Sports, Entertainment: outcomes driven by sentiment and private signals).
We proxy adverse selection by the post-trade price velocity $|\Delta P_{t+1}|$
within each market, and by the lag-1 autocorrelation of signed price changes
(negative autocorrelation $\Rightarrow$ reversal $\Rightarrow$ informed-flow pricing).

\begin{figure}[H]
  \centerline{\includegraphics[width=\columnwidth]{output/cp11_category_adverse_selection/kalshi/kalshi_cp11_category_adverse_selection.pdf}}
  \caption{Post-trade price velocity and lag-1 autocorrelation by market category.
           All categories exhibit negative lag-1 autocorrelation (bid-ask bounce).
           Finance markets show higher post-trade price velocity than Sports markets,
           consistent with more information-sensitive order flow in scheduled events.}
  \label{fig:adverse}
\end{figure}

\subsection{Sports Sentiment Arbitrage}
\label{sec:sports_sentiment}

\begin{figure*}[!t]
  \centerline{\includegraphics[width=\textwidth]{output/cp13_sports_sentiment_arbitrage/kalshi/kalshi_cp13_sports_sentiment_arbitrage.pdf}}
  \caption{Sports sentiment arbitrage decomposition across ten sport types.
           \textit{Top}: Maker (green) and taker (blue) VWR by sport, sorted by volume;
           contract volume in billions annotated at right.
           \textit{Bottom-left}: Maker VWR by time-to-resolution bucket with OLS trend line;
           the flat slope ($p = 0.88$) rejects the market-learning hypothesis.
           \textit{Bottom-right}: YES capital share by sport; all sports exceed
           the 50\% neutral line, confirming universal fan-enthusiasm bias.}
  \label{fig:sports_sentiment}
\end{figure*}

Sports markets account for approximately 74\% of all Kalshi contract volume
(12.8~billion contracts), making them the dominant driver of the platform-wide
wealth transfer. We decompose sports into ten sport types to test whether
fan-driven sentiment creates exploitable, sport-specific mispricings.
We term this mechanism \emph{sentiment arbitrage} \cite{thaler1988,camerer1998}.

\textbf{Cross-sport alpha variation.} Figure~\ref{fig:sports_sentiment} shows
that maker alpha varies substantially across sports. UFC/Boxing exhibits the
highest maker VWR at $+3.43$\,pp, consistent with the intense, binary, and
emotionally charged nature of combat sports betting. Soccer ($+1.39$\,pp) and
NFL ($+1.27$\,pp) follow, while MLB shows the lowest maker alpha ($+0.18$\,pp).
The MLB result is consistent with a relatively sophisticated sports-betting
community with deep statistical analysis traditions, whereas combat sports
attract more emotionally motivated retail flow.

\textbf{Universal YES-side fan bias.} A striking finding is that \emph{all}
sports exhibit YES capital shares well above the 50\% neutral benchmark and
above the cross-category average of 63.8\%. UFC/Boxing registers the highest
YES share at 81.8\%, followed by NHL (76.7\%) and Soccer (76.2\%). Even the
most ``efficient'' sport (MLB) shows a YES share of 74.5\%. This universal fan-enthusiasm bias provides the demand imbalance that
professional makers exploit by resting NO-side limit orders: bettors
systematically choose the affirmative (win) side regardless of the
contract's implied probability.

\textbf{Sentiment persistence.} We test whether maker alpha decays as events
approach resolution, the prediction of a market-learning hypothesis.
The time-to-resolution OLS slope across the six time buckets is
$-0.021$\,pp/bin ($p = 0.88$, $R^2 = 0.006$), statistically indistinguishable
from zero. Maker alpha is approximately constant from 30 days out through the
final six hours before resolution. This finding rejects the market-learning
hypothesis: sports bettors do not become measurably more rational as game-day
approaches, and professional makers earn a persistent structural premium
throughout the contract lifecycle.

% ---------------------------------------------------------------
%  X. LIMITATIONS
% ---------------------------------------------------------------
\section{Limitations}
\label{sec:limitations}

\textbf{Survival bias.} PnL analyses are restricted to finalized markets
(result $\in$ \{yes, no\}). Cancelled or voided markets are excluded,
potentially biasing maker alpha estimates upward.

\textbf{Fee omission.} Kalshi's taker fee schedule is not reflected in the
raw \texttt{yes\_price} field. All return calculations are gross of fees;
net-of-fee taker returns are more negative than reported.

\textbf{Position limits.} Kalshi imposes per-participant position limits.
Large, informed bets may be censored at the margin, understating the true
magnitude of adverse selection costs.

\textbf{Maker identity.} The public dataset does not contain maker identifiers,
precluding a true provider-level HHI. Our concentration metric is computed
at the contract (ticker) level as a proxy.

\textbf{Cross-platform arbitrage.} We do not model arbitrage linkages between
Kalshi and Polymarket, which may partially eliminate persistent mispricing
in practice.

\textbf{Polymarket structural differences.} Polymarket operates a hybrid
AMM/CLOB system; the maker-taker framework does not directly apply to legacy
FPMM trades, limiting cross-platform PnL comparisons.

% ---------------------------------------------------------------
%  XI. CONCLUSION
% ---------------------------------------------------------------
\section{Conclusion}
\label{sec:conclusion}

This paper characterises the microstructure of binary prediction markets using
\numTradesMillions{} million trades from Kalshi (July 2021 through November 2025)
and 404.5 million trades from Polymarket. Six robust findings emerge:
(i) a Favorite-Longshot Bias that systematically overprices low-probability
contracts, with Polymarket additionally showing favorite underpricing above 55\%;
(ii) a \$142.5M structural maker-taker wealth transfer exceeding bid-ask spread
compensation by a statistically significant margin;
(iii) cross-sectional heterogeneity in maker alpha consistent with adverse
selection, ranging from $+0.29$\,pp in Finance to $+3.01$\,pp in World Events;
(iv) a 63.8\% YES capital tilt in taker order flow, with YES-side takers earning
significantly lower returns than NO-side takers;
(v) a statistically significant upward trend in maker alpha
($\hat{\gamma}_1 = +8.21$\,bp/month, $R^2 = 0.70$), inconsistent with
efficiency-convergence predictions; and
(vi) sport-specific sentiment arbitrage in sports markets (74\% of volume),
with maker alpha ranging from $+0.18$\,pp (MLB) to $+3.43$\,pp (UFC/Boxing)
and no significant time-to-resolution decay.

Taken together, these results indicate that binary prediction markets remain
structurally biased against uninformed retail participants. Platform designers
and regulators should consider mechanisms to improve price calibration,
reduce YES-side framing effects, and increase transparency around the
structural advantages of professional liquidity providers.

% ---------------------------------------------------------------
%  ACKNOWLEDGMENT
% ---------------------------------------------------------------
\section*{Acknowledgment}

The authors thank the India Finance Conference organizing committee for the
opportunity to present this work. Trade data from Kalshi were obtained via
the Kalshi public API; Polymarket on-chain trade records were sourced from
the Polygon blockchain. Computational analysis used DuckDB, Python, and
the pandas/matplotlib/scipy open-source stack. All errors are our own.

% ---------------------------------------------------------------
%  REFERENCES
% ---------------------------------------------------------------
\begin{thebibliography}{00}

\bibitem{ali1977} F.~D.~Ali, ``Probability and utility estimates for racetrack bettors,''
  \emph{Journal of Political Economy}, vol.~85, no.~4, pp.~803--815, 1977.

\bibitem{thaler1988} R.~H.~Thaler and W.~T.~Ziemba, ``Parimutuel betting markets:
  Racetracks and lotteries,'' \emph{Journal of Economic Perspectives}, vol.~2,
  no.~2, pp.~161--174, 1988.

\bibitem{wolfers2004} J.~Wolfers and E.~Zitzewitz, ``Prediction markets,''
  \emph{Journal of Economic Perspectives}, vol.~18, no.~2, pp.~107--126, 2004.

\bibitem{arrow2008} K.~J.~Arrow et al., ``The promise of prediction markets,''
  \emph{Science}, vol.~320, pp.~877--878, 2008.

\bibitem{glosten1985} L.~R.~Glosten and P.~R.~Milgrom, ``Bid, ask and transaction
  prices in a specialist market with heterogeneously informed traders,''
  \emph{Journal of Financial Economics}, vol.~14, no.~1, pp.~71--100, 1985.

\bibitem{kyle1985} A.~S.~Kyle, ``Continuous auctions and insider trading,''
  \emph{Econometrica}, vol.~53, no.~6, pp.~1315--1335, 1985.

\bibitem{snowberg2010} E.~Snowberg and J.~Wolfers, ``Explaining the
  favorite-longshot bias: Is it risk-love or misweighting?''
  \emph{Journal of Political Economy}, vol.~118, no.~4, pp.~723--746, 2010.

\bibitem{manski2006} C.~F.~Manski, ``Interpreting the predictions of prediction
  markets,'' \emph{Economics Letters}, vol.~91, no.~3, pp.~425--429, 2006.

\bibitem{kahneman1979} D.~Kahneman and A.~Tversky, ``Prospect theory: An analysis
  of decision under risk,'' \emph{Econometrica}, vol.~47, no.~2,
  pp.~263--292, 1979.

\bibitem{ottaviani2008} M.~Ottaviani and P.~N.~Sorensen, ``The favorite-longshot
  bias: An overview of the main explanations,'' \emph{Handbook of Sports and
  Lottery Markets}, 2008.

\bibitem{amihud2002} Y.~Amihud, ``Illiquidity and stock returns: Cross-section
  and time-series effects,'' \emph{Journal of Financial Markets}, vol.~5,
  no.~1, pp.~31--56, 2002.

\bibitem{roll1984} R.~Roll, ``A simple implicit measure of the effective
  bid-ask spread in an efficient market,'' \emph{Journal of Finance},
  vol.~39, no.~4, pp.~1127--1139, 1984.

\bibitem{fama1970} E.~F.~Fama, ``Efficient capital markets: A review of theory
  and empirical work,'' \emph{Journal of Finance}, vol.~25, no.~2,
  pp.~383--417, 1970.

\bibitem{ohara1995} M.~O'Hara, \emph{Market Microstructure Theory}.
  Blackwell, Cambridge, MA, 1995.

\bibitem{camerer1998} C.~F.~Camerer, ``Can asset markets be manipulated?
  A field experiment with racetrack betting,'' \emph{Journal of Political
  Economy}, vol.~106, no.~3, pp.~457--482, 1998.

\bibitem{berg2008} J.~Berg, F.~Nelson, and T.~Rietz, ``Prediction market
  accuracy in the long run,'' \emph{International Journal of Forecasting},
  vol.~24, no.~2, pp.~285--300, 2008.

\bibitem{hanson2003} R.~Hanson, ``Combinatorial information market design,''
  \emph{Information Systems Frontiers}, vol.~5, no.~1, pp.~107--119, 2003.

\bibitem{hong2007} H.~Hong and J.~C.~Stein, ``Disagreement and the stock market,''
  \emph{Journal of Economic Perspectives}, vol.~21, no.~2, pp.~109--128, 2007.

\end{thebibliography}

\end{document}
